\documentclass[11pt,a4paper]{article}
% \usepackage[utf8]{inputenc}
\usepackage{geometry}
\usepackage{graphicx}
\usepackage{booktabs}
\usepackage{hyperref}
\usepackage{amsmath}
\usepackage{amssymb}
\usepackage{natbib}
\usepackage{enumitem}
\usepackage{xcolor}
\usepackage{fancyhdr}
\usepackage{titlesec}

\usepackage{fontspec}
\usepackage{xeCJK}

\setmainfont{Times New Roman}

% ================= 图像使用规范（供生成模型参考） =================
% 使用如下标准格式插入图片：
%
% \begin{figure}[htbp]
% \centering
% \includegraphics[width=0.8\linewidth]{path/to/image}
% \caption{对图像进行描述，并解释其在文献综述中的作用}
% \label{fig:unique_label}
% \end{figure}
%
% 要求：
% 1. 图片路径必须使用原始路径（不要修改或重写）
% 2. 每个图片必须有 \caption 和 \label
% 3. 在正文中使用如下方式引用：
%    Figure~\ref{fig:unique_label}
% 4. 图片应放置在相关内容附近（不要集中在文末）
% ============================================================

% 页面设置
\geometry{left=2.5cm,right=2.5cm,top=2.5cm,bottom=2.5cm}

% 页眉页脚
\pagestyle{fancy}
\fancyhf{}
\fancyhead[L]{\leftmark}
\fancyhead[R]{文献综述}
\fancyfoot[C]{\thepage}
\renewcommand{\headrulewidth}{0.4pt}

% 标题格式
\titleformat{\section}{\Large\bfseries\color{blue!50!black}}{\thesection}{1em}{}
\titleformat{\subsection}{\large\bfseries\color{blue!30!black}}{\thesubsection}{1em}{}

% 超链接颜色
\hypersetup{
    colorlinks=true,
    linkcolor=blue!60!black,
    citecolor=green!50!black,
    urlcolor=blue!70!black
}

\begin{filecontents}{references.bib}
% 示例参考文献，请替换为实际引用
@article{example2024,
  title={用于文献综述的示例论文},
  author={Author, Example},
  journal={示例期刊},
  volume={1},
  number={1},
  pages={1--10},
  year={2024},
  publisher={示例出版社}
}

@inproceedings{sample2023,
  title={示例会议论文},
  author={Sample, Author and Test, Researcher},
  booktitle={示例会议论文集},
  pages={100--110},
  year={2023}
}
\end{filecontents}

% ========== 标题信息 ==========
\title{\textbf{文献综述：[你的研究主题]}}
\author{你的姓名 \\
        你的机构 \\
        \texttt{your.email@example.com}}
\date{\today}

\begin{document}

\maketitle

% ========== 摘要 ==========
\begin{abstract}
\noindent
\textbf{研究目的：} 简要描述本综述的目标与范围。 \\
\textbf{研究方法：} 描述文献检索策略与筛选标准。 \\
\textbf{主要发现：} 总结文献中的关键研究结论。 \\
\textbf{结论：} 给出主要结论及其意义。

\vspace{0.5em}
\noindent\textbf{关键词：} 关键词1，关键词2，关键词3，关键词4
\end{abstract}

\tableofcontents
\newpage

% ========== 1. 引言 ==========
\section{引言}
\label{sec:intro}

\subsection{研究背景}
介绍研究主题的背景，说明其重要性及在相关领域中的意义。

\subsection{研究问题}
明确本综述试图回答的问题：
\begin{enumerate}
    \item 主要研究问题是什么？
    \item 聚焦该主题的哪些具体方面？
    \item 试图识别哪些研究空白？
\end{enumerate}

\subsection{范围与局限}
界定综述范围：
\begin{itemize}
    \item 时间范围（例如：2010–2024）
    \item 文献类型（期刊、会议论文等）
    \item 地理或领域限制
    \item 排除标准
\end{itemize}

% ========== 2. 研究方法 ==========
\section{研究方法}
\label{sec:method}

\subsection{检索策略}
描述系统性检索方法：
\begin{itemize}
    \item 使用的数据库（如 Google Scholar、IEEE Xplore、PubMed）
    \item 检索关键词
    \item 检索时间
    \item 是否采用滚雪球或引用追踪
\end{itemize}

\subsection{筛选标准}
说明纳入与排除标准：
\begin{itemize}
    \item \textbf{纳入标准：} 文献需满足的条件
    \item \textbf{排除标准：} 文献被剔除的条件
\end{itemize}

\subsection{数据提取}
说明如何从文献中提取并组织信息。

% ========== 3. 文献分类框架 ==========
\section{分类框架}
\label{sec:framework}

介绍组织文献的分类方法：

\subsection{按时间}
按时间发展阶段进行组织。

\subsection{按方法}
按研究方法分类（实验、理论、综述等）。

\subsection{按主题}
按研究主题或问题分类。

% ========== 4. 文献综述 ==========
\section{文献综述}
\label{sec:review}

\subsection{主题1：[第一个主题]}
总结该主题的关键文献：
\begin{itemize}
    \item 研究目标
    \item 方法
    \item 主要贡献
    \item 优势与不足
\end{itemize}

\textbf{关键文献：}
\begin{itemize}
    \item \cite{example2024} - 简要说明贡献
    \item \cite{sample2023} - 简要说明贡献
\end{itemize}

\subsection{主题2：[第二个主题]}
（结构同上）

\subsection{主题3：[第三个主题]}

% ========== 5. 比较分析 ==========
\section{比较分析}
\label{sec:comparison}

\subsection{跨主题比较}
\begin{table}[htbp]
\centering
\caption{方法比较}
\label{tab:comparison}
\begin{tabular}{@{}llll@{}}
\toprule
方法 & 优点 & 缺点 & 适用场景 \\
\midrule
方法A & 优点1 & 缺点1 & 场景X \\
方法B & 优点2 & 缺点2 & 场景Y \\
方法C & 优点3 & 缺点3 & 场景Z \\
\bottomrule
\end{tabular}
\end{table}

\subsection{趋势分析}
\begin{itemize}
    \item 新兴方向
    \item 方法演化
    \item 研究热度变化
\end{itemize}

% ========== 6. 研究空白 ==========
\section{研究空白}
\label{sec:gaps}

\begin{enumerate}
    \item \textbf{空白1：} 描述
    \item \textbf{空白2：} 描述
    \item \textbf{空白3：} 描述
\end{enumerate}

% ========== 7. 讨论 ==========
\section{讨论}
\label{sec:discussion}

\subsection{综合分析}
总结整体结论。

\subsection{理论意义}
讨论理论贡献。

\subsection{实践意义}
讨论实际应用价值。

% ========== 8. 结论 ==========
\section{结论}
\label{sec:conclusion}

\subsection{总结}
\begin{itemize}
    \item 主要主题
    \item 关键发现
    \item 研究空白
\end{itemize}

\subsection{未来方向}
\begin{enumerate}
    \item 重要研究问题
    \item 方法改进
    \item 新兴方向
\end{enumerate}

\subsection{局限性}
说明综述本身的局限。

% ========== 参考文献 ==========
\newpage
\bibliographystyle{apalike}
\bibliography{references}

% ========== 附录 ==========
\newpage
\appendix
\section{检索结果}
\label{app:search}

\section{数据提取表}
\label{app:data}

\end{document}